% Include after subsection 5.8, immediately before Section 6:
% % Include after subsection 5.8, immediately before Section 6:
% % Include after subsection 5.8, immediately before Section 6:
% % Include after subsection 5.8, immediately before Section 6:
% \input{implications_of_marginality}
% Uses the parent paper's marginalbox, boxblue, boxlightblue, and TikZ setup.
% All equations and the analytic figure are unnumbered.
% Starts a new page and uses one column for each subsubsection.
\clearpage
\subsection{Implications of marginality}
\label{sxn:implications_marginality}

\subsubsection{The model-average exponent}
\label{sxn:model_average_marginality}
SETOL combines positive layer qualities into the model-quality approximation
$\overline Q_{\rm NN}\simeq\prod_{\ell=1}^K\overline Q_\ell$
[Eq.~\eqref{eq:setol_layer_product}]. Its logarithm converts the
layer product into a sum, giving the average logarithmic quality,
\[
 \frac{1}{K}\log\overline Q_{\rm NN}
 \simeq\frac{1}{K}\sum_{\ell=1}^K\log\overline Q_\ell .
\]
The layer product motivates averaging the logarithmic gains of
the spectral bands. At a fixed checkpoint, Eq.~\eqref{eq:shellpower} gives
$g_{E,\ell}(\Lambda_\ell)\propto\Lambda_\ell^{2-\alpha_\ell}$.
Comparing $\Lambda_\ell$ with $b\Lambda_\ell$ inside each fitted window
multiplies the gain by $b^{2-\alpha_\ell}$. A common factor
$b>1$ gives every layer the same step in logarithmic scale.
The powers of their product then add.

\begin{tcolorbox}[marginalbox,title=Marginality on average]
\[
\begin{gathered}
 \bar\alpha:=\frac{1}{K}\sum_{\ell=1}^K\alpha_\ell,\qquad
 \prod_{\ell=1}^K b^{2-\alpha_\ell}=b^{K(2-\bar\alpha)},\\[3pt]
 \bar y_E:=\frac{1}{K}\sum_{\ell=1}^K(2-\alpha_\ell)
          =2-\bar\alpha .
\end{gathered}
\]
\textcolor{boxblue}{The product of the band gains is marginal
when $\bar\alpha=2$. The mean logarithmic band slope then vanishes.}
\end{tcolorbox}

This average expresses model marginality as a balance across layers.
Layers with $\alpha_\ell=2$ give $\bar\alpha=2$. Equal and opposite
deviations from two also yield this average through compensation.

Two layers with $\alpha_1=1.5$ and $\alpha_2=2.5$ achieve
$\bar\alpha=2$ through compensation. Doubling the spectral scale
multiplies their band gains by $\sqrt{2}$ and $1/\sqrt{2}$.
These reciprocal factors have product one and balance the layer slopes.

\begin{center}
\resizebox{.96\columnwidth}{!}{%
\begin{tikzpicture}[x=1cm,y=1cm,font=\sffamily\small,>=Stealth]
 \path[use as bounding box] (-.85,-.62) rectangle (7.0,3.28);
 \node[text=boxblue,font=\sffamily\bfseries\small] at (3.05,3.0)
   {Layers balance their spectral slopes};
 \draw[->,gray] (0,0) -- (6.7,0);
 \draw[->,gray] (0,0) -- (0,2.65);
 \foreach \x/\lab in {0/1,2.4/2,4.8/4}{
   \draw[gray] (\x,0) -- (\x,-.08);
   \node[below,font=\small] at (\x,-.08) {$\lab$};
 }
 \foreach \y/\lab in {.3/{1/2},1.3/{1},2.3/{2}}{
   \draw[gray] (0,\y) -- (-.08,\y);
   \node[left,font=\small] at (-.08,\y) {$\lab$};
 }
 \draw[boxblue,line width=1.2pt] (0,1.3) -- (4.8,2.3);
 \draw[orange!80!black,line width=1.2pt] (0,1.3) -- (4.8,.3);
 \draw[black,dashed,line width=.9pt] (0,1.3) -- (4.8,1.3);
 \node[right,text=boxblue] at (4.86,2.3) {$\alpha_1=1.5$};
 \node[right,text=orange!80!black] at (4.86,.3) {$\alpha_2=2.5$};
 \node[right,font=\sffamily\scriptsize] at (4.86,1.3) {geometric mean};
 \node[below,font=\sffamily\scriptsize] at (3.0,-.35)
   {spectral scale factor $b$ (log scale)};
\end{tikzpicture}}
\captionsetup{justification=raggedright,singlelinecheck=false}
\captionof*{figure}{\textbf{Marginal on average.}
The band-gain ratios $g_{E,\ell}(b\Lambda_\ell)/g_{E,\ell}(\Lambda_\ell)$
have a geometric mean of one. Both axes use logarithmic scales.}
\end{center}

The mean exponent diagnoses average marginality, and its spread
describes how the layers achieve that balance. Exponents clustered
near two indicate coordinated layerwise marginality. Exponents
on opposite sides of two indicate compensation.

\newpage
\subsubsection{Mode seeking versus geometry seeking}
\label{sxn:mean_variance_marginality}
The spectral moments place the mean and variance at different
marginal boundaries. For a tail
$\rho(\lambda)\simeq C\lambda^{-\alpha}$ with $\alpha>1$,
we extend the upper cutoff $\Lambda$ while holding the lower cutoff
$\lambda_0>0$ and amplitude $C$ fixed. The tail contribution to the
$k$th raw moment obeys
\[
\begin{aligned}
 \mu_k(\Lambda)&:=\int_{\lambda_0}^{\Lambda}
                  \lambda^k\rho(\lambda)\,d\lambda,\\
 \frac{d\mu_k}{d\log\Lambda}
   &=\Lambda^{k+1}\rho(\Lambda)
     \simeq C\Lambda^{k+1-\alpha}.
\end{aligned}
\]
Equal contributions from successive logarithmic bands define
moment marginality. At $\alpha=k+1$, these contributions give
$\mu_k(\Lambda)\sim C\log(\Lambda/\lambda_0)$.
For $\alpha>k+1$, the moment converges. Thus $\alpha=2$ is marginal for the first moment and
$\alpha=3$ is marginal for the second moment. Since the mean is
finite at $\alpha=3$, the variance inherits the second moment's
logarithmic growth.
Both ideal moments are finite for $\alpha>3$.

\paragraph{Mode seeking.}
We hypothesize that layers with $1<\alpha<2$ seek particular
modes of the training data and become more prone to overfitting.
Their upper spectral bands contribute progressively more
first-moment energy, favoring concentration of squared gain in a
few leading singular directions. These
directions amplify training-specific patterns and concentrate the
layer's response on them. At $\alpha=2$, equal logarithmic bands
carry equal first-moment energy, balancing that gain across
spectral scales.

\paragraph{Geometry seeking.}
We hypothesize that layers with $2<\alpha<3$ seek the geometry of
the data by amplifying structured variation across directions.
Their ideal mean is finite while their variance grows with the
upper cutoff. Higher bands contribute less to
the mean but more to the second moment. At $\alpha=2.5$, for
example, moving up one octave multiplies the first-moment band
weight by $1/\sqrt{2}$ and its second-moment weight by $\sqrt{2}$.

This spread measures how unevenly a layer amplifies its singular
directions. For the full
mean-normalized correlation operator $\widetilde{\mathbf X}$
with $M$ eigenvalues and $\langle\tilde\lambda\rangle=1$,
\[
 \operatorname{Var}_{\rm ESD}(\tilde\lambda)
   =\frac{1}{M}\sum_{i=1}^M(\tilde\lambda_i-1)^2
   =\frac{1}{M}\|\widetilde{\mathbf X}-I_M\|_F^2 .
\]
The normalized spectrum expresses these hypotheses through
concentration and directional spread. The fraction of the trace
carried by the leading modes measures concentration. The variance
of the eigenvalues measures the spread of squared directional
gains. In the proposed learning picture, mode seeking emphasizes
the first quantity and geometry seeking emphasizes the second.
\clearpage

% Uses the parent paper's marginalbox, boxblue, boxlightblue, and TikZ setup.
% All equations and the analytic figure are unnumbered.
% Starts a new page and uses one column for each subsubsection.
\clearpage
\subsection{Implications of marginality}
\label{sxn:implications_marginality}

\subsubsection{The model-average exponent}
\label{sxn:model_average_marginality}
SETOL combines positive layer qualities into the model-quality approximation
$\overline Q_{\rm NN}\simeq\prod_{\ell=1}^K\overline Q_\ell$
[Eq.~\eqref{eq:setol_layer_product}]. Its logarithm converts the
layer product into a sum, giving the average logarithmic quality,
\[
 \frac{1}{K}\log\overline Q_{\rm NN}
 \simeq\frac{1}{K}\sum_{\ell=1}^K\log\overline Q_\ell .
\]
The layer product motivates averaging the logarithmic gains of
the spectral bands. At a fixed checkpoint, Eq.~\eqref{eq:shellpower} gives
$g_{E,\ell}(\Lambda_\ell)\propto\Lambda_\ell^{2-\alpha_\ell}$.
Comparing $\Lambda_\ell$ with $b\Lambda_\ell$ inside each fitted window
multiplies the gain by $b^{2-\alpha_\ell}$. A common factor
$b>1$ gives every layer the same step in logarithmic scale.
The powers of their product then add.

\begin{tcolorbox}[marginalbox,title=Marginality on average]
\[
\begin{gathered}
 \bar\alpha:=\frac{1}{K}\sum_{\ell=1}^K\alpha_\ell,\qquad
 \prod_{\ell=1}^K b^{2-\alpha_\ell}=b^{K(2-\bar\alpha)},\\[3pt]
 \bar y_E:=\frac{1}{K}\sum_{\ell=1}^K(2-\alpha_\ell)
          =2-\bar\alpha .
\end{gathered}
\]
\textcolor{boxblue}{The product of the band gains is marginal
when $\bar\alpha=2$. The mean logarithmic band slope then vanishes.}
\end{tcolorbox}

This average expresses model marginality as a balance across layers.
Layers with $\alpha_\ell=2$ give $\bar\alpha=2$. Equal and opposite
deviations from two also yield this average through compensation.

Two layers with $\alpha_1=1.5$ and $\alpha_2=2.5$ achieve
$\bar\alpha=2$ through compensation. Doubling the spectral scale
multiplies their band gains by $\sqrt{2}$ and $1/\sqrt{2}$.
These reciprocal factors have product one and balance the layer slopes.

\begin{center}
\resizebox{.96\columnwidth}{!}{%
\begin{tikzpicture}[x=1cm,y=1cm,font=\sffamily\small,>=Stealth]
 \path[use as bounding box] (-.85,-.62) rectangle (7.0,3.28);
 \node[text=boxblue,font=\sffamily\bfseries\small] at (3.05,3.0)
   {Layers balance their spectral slopes};
 \draw[->,gray] (0,0) -- (6.7,0);
 \draw[->,gray] (0,0) -- (0,2.65);
 \foreach \x/\lab in {0/1,2.4/2,4.8/4}{
   \draw[gray] (\x,0) -- (\x,-.08);
   \node[below,font=\small] at (\x,-.08) {$\lab$};
 }
 \foreach \y/\lab in {.3/{1/2},1.3/{1},2.3/{2}}{
   \draw[gray] (0,\y) -- (-.08,\y);
   \node[left,font=\small] at (-.08,\y) {$\lab$};
 }
 \draw[boxblue,line width=1.2pt] (0,1.3) -- (4.8,2.3);
 \draw[orange!80!black,line width=1.2pt] (0,1.3) -- (4.8,.3);
 \draw[black,dashed,line width=.9pt] (0,1.3) -- (4.8,1.3);
 \node[right,text=boxblue] at (4.86,2.3) {$\alpha_1=1.5$};
 \node[right,text=orange!80!black] at (4.86,.3) {$\alpha_2=2.5$};
 \node[right,font=\sffamily\scriptsize] at (4.86,1.3) {geometric mean};
 \node[below,font=\sffamily\scriptsize] at (3.0,-.35)
   {spectral scale factor $b$ (log scale)};
\end{tikzpicture}}
\captionsetup{justification=raggedright,singlelinecheck=false}
\captionof*{figure}{\textbf{Marginal on average.}
The band-gain ratios $g_{E,\ell}(b\Lambda_\ell)/g_{E,\ell}(\Lambda_\ell)$
have a geometric mean of one. Both axes use logarithmic scales.}
\end{center}

The mean exponent diagnoses average marginality, and its spread
describes how the layers achieve that balance. Exponents clustered
near two indicate coordinated layerwise marginality. Exponents
on opposite sides of two indicate compensation.

\newpage
\subsubsection{Mode seeking versus geometry seeking}
\label{sxn:mean_variance_marginality}
The spectral moments place the mean and variance at different
marginal boundaries. For a tail
$\rho(\lambda)\simeq C\lambda^{-\alpha}$ with $\alpha>1$,
we extend the upper cutoff $\Lambda$ while holding the lower cutoff
$\lambda_0>0$ and amplitude $C$ fixed. The tail contribution to the
$k$th raw moment obeys
\[
\begin{aligned}
 \mu_k(\Lambda)&:=\int_{\lambda_0}^{\Lambda}
                  \lambda^k\rho(\lambda)\,d\lambda,\\
 \frac{d\mu_k}{d\log\Lambda}
   &=\Lambda^{k+1}\rho(\Lambda)
     \simeq C\Lambda^{k+1-\alpha}.
\end{aligned}
\]
Equal contributions from successive logarithmic bands define
moment marginality. At $\alpha=k+1$, these contributions give
$\mu_k(\Lambda)\sim C\log(\Lambda/\lambda_0)$.
For $\alpha>k+1$, the moment converges. Thus $\alpha=2$ is marginal for the first moment and
$\alpha=3$ is marginal for the second moment. Since the mean is
finite at $\alpha=3$, the variance inherits the second moment's
logarithmic growth.
Both ideal moments are finite for $\alpha>3$.

\paragraph{Mode seeking.}
We hypothesize that layers with $1<\alpha<2$ seek particular
modes of the training data and become more prone to overfitting.
Their upper spectral bands contribute progressively more
first-moment energy, favoring concentration of squared gain in a
few leading singular directions. These
directions amplify training-specific patterns and concentrate the
layer's response on them. At $\alpha=2$, equal logarithmic bands
carry equal first-moment energy, balancing that gain across
spectral scales.

\paragraph{Geometry seeking.}
We hypothesize that layers with $2<\alpha<3$ seek the geometry of
the data by amplifying structured variation across directions.
Their ideal mean is finite while their variance grows with the
upper cutoff. Higher bands contribute less to
the mean but more to the second moment. At $\alpha=2.5$, for
example, moving up one octave multiplies the first-moment band
weight by $1/\sqrt{2}$ and its second-moment weight by $\sqrt{2}$.

This spread measures how unevenly a layer amplifies its singular
directions. For the full
mean-normalized correlation operator $\widetilde{\mathbf X}$
with $M$ eigenvalues and $\langle\tilde\lambda\rangle=1$,
\[
 \operatorname{Var}_{\rm ESD}(\tilde\lambda)
   =\frac{1}{M}\sum_{i=1}^M(\tilde\lambda_i-1)^2
   =\frac{1}{M}\|\widetilde{\mathbf X}-I_M\|_F^2 .
\]
The normalized spectrum expresses these hypotheses through
concentration and directional spread. The fraction of the trace
carried by the leading modes measures concentration. The variance
of the eigenvalues measures the spread of squared directional
gains. In the proposed learning picture, mode seeking emphasizes
the first quantity and geometry seeking emphasizes the second.
\clearpage

% Uses the parent paper's marginalbox, boxblue, boxlightblue, and TikZ setup.
% All equations and the analytic figure are unnumbered.
% Starts a new page and uses one column for each subsubsection.
\clearpage
\subsection{Implications of marginality}
\label{sxn:implications_marginality}

\subsubsection{The model-average exponent}
\label{sxn:model_average_marginality}
SETOL combines positive layer qualities into the model-quality approximation
$\overline Q_{\rm NN}\simeq\prod_{\ell=1}^K\overline Q_\ell$
[Eq.~\eqref{eq:setol_layer_product}]. Its logarithm converts the
layer product into a sum, giving the average logarithmic quality,
\[
 \frac{1}{K}\log\overline Q_{\rm NN}
 \simeq\frac{1}{K}\sum_{\ell=1}^K\log\overline Q_\ell .
\]
The layer product motivates averaging the logarithmic gains of
the spectral bands. At a fixed checkpoint, Eq.~\eqref{eq:shellpower} gives
$g_{E,\ell}(\Lambda_\ell)\propto\Lambda_\ell^{2-\alpha_\ell}$.
Comparing $\Lambda_\ell$ with $b\Lambda_\ell$ inside each fitted window
multiplies the gain by $b^{2-\alpha_\ell}$. A common factor
$b>1$ gives every layer the same step in logarithmic scale.
The powers of their product then add.

\begin{tcolorbox}[marginalbox,title=Marginality on average]
\[
\begin{gathered}
 \bar\alpha:=\frac{1}{K}\sum_{\ell=1}^K\alpha_\ell,\qquad
 \prod_{\ell=1}^K b^{2-\alpha_\ell}=b^{K(2-\bar\alpha)},\\[3pt]
 \bar y_E:=\frac{1}{K}\sum_{\ell=1}^K(2-\alpha_\ell)
          =2-\bar\alpha .
\end{gathered}
\]
\textcolor{boxblue}{The product of the band gains is marginal
when $\bar\alpha=2$. The mean logarithmic band slope then vanishes.}
\end{tcolorbox}

This average expresses model marginality as a balance across layers.
Layers with $\alpha_\ell=2$ give $\bar\alpha=2$. Equal and opposite
deviations from two also yield this average through compensation.

Two layers with $\alpha_1=1.5$ and $\alpha_2=2.5$ achieve
$\bar\alpha=2$ through compensation. Doubling the spectral scale
multiplies their band gains by $\sqrt{2}$ and $1/\sqrt{2}$.
These reciprocal factors have product one and balance the layer slopes.

\begin{center}
\resizebox{.96\columnwidth}{!}{%
\begin{tikzpicture}[x=1cm,y=1cm,font=\sffamily\small,>=Stealth]
 \path[use as bounding box] (-.85,-.62) rectangle (7.0,3.28);
 \node[text=boxblue,font=\sffamily\bfseries\small] at (3.05,3.0)
   {Layers balance their spectral slopes};
 \draw[->,gray] (0,0) -- (6.7,0);
 \draw[->,gray] (0,0) -- (0,2.65);
 \foreach \x/\lab in {0/1,2.4/2,4.8/4}{
   \draw[gray] (\x,0) -- (\x,-.08);
   \node[below,font=\small] at (\x,-.08) {$\lab$};
 }
 \foreach \y/\lab in {.3/{1/2},1.3/{1},2.3/{2}}{
   \draw[gray] (0,\y) -- (-.08,\y);
   \node[left,font=\small] at (-.08,\y) {$\lab$};
 }
 \draw[boxblue,line width=1.2pt] (0,1.3) -- (4.8,2.3);
 \draw[orange!80!black,line width=1.2pt] (0,1.3) -- (4.8,.3);
 \draw[black,dashed,line width=.9pt] (0,1.3) -- (4.8,1.3);
 \node[right,text=boxblue] at (4.86,2.3) {$\alpha_1=1.5$};
 \node[right,text=orange!80!black] at (4.86,.3) {$\alpha_2=2.5$};
 \node[right,font=\sffamily\scriptsize] at (4.86,1.3) {geometric mean};
 \node[below,font=\sffamily\scriptsize] at (3.0,-.35)
   {spectral scale factor $b$ (log scale)};
\end{tikzpicture}}
\captionsetup{justification=raggedright,singlelinecheck=false}
\captionof*{figure}{\textbf{Marginal on average.}
The band-gain ratios $g_{E,\ell}(b\Lambda_\ell)/g_{E,\ell}(\Lambda_\ell)$
have a geometric mean of one. Both axes use logarithmic scales.}
\end{center}

The mean exponent diagnoses average marginality, and its spread
describes how the layers achieve that balance. Exponents clustered
near two indicate coordinated layerwise marginality. Exponents
on opposite sides of two indicate compensation.

\newpage
\subsubsection{Mode seeking versus geometry seeking}
\label{sxn:mean_variance_marginality}
The spectral moments place the mean and variance at different
marginal boundaries. For a tail
$\rho(\lambda)\simeq C\lambda^{-\alpha}$ with $\alpha>1$,
we extend the upper cutoff $\Lambda$ while holding the lower cutoff
$\lambda_0>0$ and amplitude $C$ fixed. The tail contribution to the
$k$th raw moment obeys
\[
\begin{aligned}
 \mu_k(\Lambda)&:=\int_{\lambda_0}^{\Lambda}
                  \lambda^k\rho(\lambda)\,d\lambda,\\
 \frac{d\mu_k}{d\log\Lambda}
   &=\Lambda^{k+1}\rho(\Lambda)
     \simeq C\Lambda^{k+1-\alpha}.
\end{aligned}
\]
Equal contributions from successive logarithmic bands define
moment marginality. At $\alpha=k+1$, these contributions give
$\mu_k(\Lambda)\sim C\log(\Lambda/\lambda_0)$.
For $\alpha>k+1$, the moment converges. Thus $\alpha=2$ is marginal for the first moment and
$\alpha=3$ is marginal for the second moment. Since the mean is
finite at $\alpha=3$, the variance inherits the second moment's
logarithmic growth.
Both ideal moments are finite for $\alpha>3$.

\paragraph{Mode seeking.}
We hypothesize that layers with $1<\alpha<2$ seek particular
modes of the training data and become more prone to overfitting.
Their upper spectral bands contribute progressively more
first-moment energy, favoring concentration of squared gain in a
few leading singular directions. These
directions amplify training-specific patterns and concentrate the
layer's response on them. At $\alpha=2$, equal logarithmic bands
carry equal first-moment energy, balancing that gain across
spectral scales.

\paragraph{Geometry seeking.}
We hypothesize that layers with $2<\alpha<3$ seek the geometry of
the data by amplifying structured variation across directions.
Their ideal mean is finite while their variance grows with the
upper cutoff. Higher bands contribute less to
the mean but more to the second moment. At $\alpha=2.5$, for
example, moving up one octave multiplies the first-moment band
weight by $1/\sqrt{2}$ and its second-moment weight by $\sqrt{2}$.

This spread measures how unevenly a layer amplifies its singular
directions. For the full
mean-normalized correlation operator $\widetilde{\mathbf X}$
with $M$ eigenvalues and $\langle\tilde\lambda\rangle=1$,
\[
 \operatorname{Var}_{\rm ESD}(\tilde\lambda)
   =\frac{1}{M}\sum_{i=1}^M(\tilde\lambda_i-1)^2
   =\frac{1}{M}\|\widetilde{\mathbf X}-I_M\|_F^2 .
\]
The normalized spectrum expresses these hypotheses through
concentration and directional spread. The fraction of the trace
carried by the leading modes measures concentration. The variance
of the eigenvalues measures the spread of squared directional
gains. In the proposed learning picture, mode seeking emphasizes
the first quantity and geometry seeking emphasizes the second.
\clearpage

% Uses the parent paper's marginalbox, boxblue, boxlightblue, and TikZ setup.
% All equations and the analytic figure are unnumbered.
% Starts a new page and uses one column for each subsubsection.
\clearpage
\subsection{Implications of marginality}
\label{sxn:implications_marginality}

\subsubsection{The model-average exponent}
\label{sxn:model_average_marginality}
SETOL combines positive layer qualities into the model-quality approximation
$\overline Q_{\rm NN}\simeq\prod_{\ell=1}^K\overline Q_\ell$
[Eq.~\eqref{eq:setol_layer_product}]. Its logarithm converts the
layer product into a sum, giving the average logarithmic quality,
\[
 \frac{1}{K}\log\overline Q_{\rm NN}
 \simeq\frac{1}{K}\sum_{\ell=1}^K\log\overline Q_\ell .
\]
The layer product motivates averaging the logarithmic gains of
the spectral bands. At a fixed checkpoint, Eq.~\eqref{eq:shellpower} gives
$g_{E,\ell}(\Lambda_\ell)\propto\Lambda_\ell^{2-\alpha_\ell}$.
Comparing $\Lambda_\ell$ with $b\Lambda_\ell$ inside each fitted window
multiplies the gain by $b^{2-\alpha_\ell}$. A common factor
$b>1$ gives every layer the same step in logarithmic scale.
The powers of their product then add.

\begin{tcolorbox}[marginalbox,title=Marginality on average]
\[
\begin{gathered}
 \bar\alpha:=\frac{1}{K}\sum_{\ell=1}^K\alpha_\ell,\qquad
 \prod_{\ell=1}^K b^{2-\alpha_\ell}=b^{K(2-\bar\alpha)},\\[3pt]
 \bar y_E:=\frac{1}{K}\sum_{\ell=1}^K(2-\alpha_\ell)
          =2-\bar\alpha .
\end{gathered}
\]
\textcolor{boxblue}{The product of the band gains is marginal
when $\bar\alpha=2$. The mean logarithmic band slope then vanishes.}
\end{tcolorbox}

This average expresses model marginality as a balance across layers.
Layers with $\alpha_\ell=2$ give $\bar\alpha=2$. Equal and opposite
deviations from two also yield this average through compensation.

Two layers with $\alpha_1=1.5$ and $\alpha_2=2.5$ achieve
$\bar\alpha=2$ through compensation. Doubling the spectral scale
multiplies their band gains by $\sqrt{2}$ and $1/\sqrt{2}$.
These reciprocal factors have product one and balance the layer slopes.

\begin{center}
\resizebox{.96\columnwidth}{!}{%
\begin{tikzpicture}[x=1cm,y=1cm,font=\sffamily\small,>=Stealth]
 \path[use as bounding box] (-.85,-.62) rectangle (7.0,3.28);
 \node[text=boxblue,font=\sffamily\bfseries\small] at (3.05,3.0)
   {Layers balance their spectral slopes};
 \draw[->,gray] (0,0) -- (6.7,0);
 \draw[->,gray] (0,0) -- (0,2.65);
 \foreach \x/\lab in {0/1,2.4/2,4.8/4}{
   \draw[gray] (\x,0) -- (\x,-.08);
   \node[below,font=\small] at (\x,-.08) {$\lab$};
 }
 \foreach \y/\lab in {.3/{1/2},1.3/{1},2.3/{2}}{
   \draw[gray] (0,\y) -- (-.08,\y);
   \node[left,font=\small] at (-.08,\y) {$\lab$};
 }
 \draw[boxblue,line width=1.2pt] (0,1.3) -- (4.8,2.3);
 \draw[orange!80!black,line width=1.2pt] (0,1.3) -- (4.8,.3);
 \draw[black,dashed,line width=.9pt] (0,1.3) -- (4.8,1.3);
 \node[right,text=boxblue] at (4.86,2.3) {$\alpha_1=1.5$};
 \node[right,text=orange!80!black] at (4.86,.3) {$\alpha_2=2.5$};
 \node[right,font=\sffamily\scriptsize] at (4.86,1.3) {geometric mean};
 \node[below,font=\sffamily\scriptsize] at (3.0,-.35)
   {spectral scale factor $b$ (log scale)};
\end{tikzpicture}}
\captionsetup{justification=raggedright,singlelinecheck=false}
\captionof*{figure}{\textbf{Marginal on average.}
The band-gain ratios $g_{E,\ell}(b\Lambda_\ell)/g_{E,\ell}(\Lambda_\ell)$
have a geometric mean of one. Both axes use logarithmic scales.}
\end{center}

The mean exponent diagnoses average marginality, and its spread
describes how the layers achieve that balance. Exponents clustered
near two indicate coordinated layerwise marginality. Exponents
on opposite sides of two indicate compensation.

\newpage
\subsubsection{Mode seeking versus geometry seeking}
\label{sxn:mean_variance_marginality}
The spectral moments place the mean and variance at different
marginal boundaries. For a tail
$\rho(\lambda)\simeq C\lambda^{-\alpha}$ with $\alpha>1$,
we extend the upper cutoff $\Lambda$ while holding the lower cutoff
$\lambda_0>0$ and amplitude $C$ fixed. The tail contribution to the
$k$th raw moment obeys
\[
\begin{aligned}
 \mu_k(\Lambda)&:=\int_{\lambda_0}^{\Lambda}
                  \lambda^k\rho(\lambda)\,d\lambda,\\
 \frac{d\mu_k}{d\log\Lambda}
   &=\Lambda^{k+1}\rho(\Lambda)
     \simeq C\Lambda^{k+1-\alpha}.
\end{aligned}
\]
Equal contributions from successive logarithmic bands define
moment marginality. At $\alpha=k+1$, these contributions give
$\mu_k(\Lambda)\sim C\log(\Lambda/\lambda_0)$.
For $\alpha>k+1$, the moment converges. Thus $\alpha=2$ is marginal for the first moment and
$\alpha=3$ is marginal for the second moment. Since the mean is
finite at $\alpha=3$, the variance inherits the second moment's
logarithmic growth.
Both ideal moments are finite for $\alpha>3$.

\paragraph{Mode seeking.}
We hypothesize that layers with $1<\alpha<2$ seek particular
modes of the training data and become more prone to overfitting.
Their upper spectral bands contribute progressively more
first-moment energy, favoring concentration of squared gain in a
few leading singular directions. These
directions amplify training-specific patterns and concentrate the
layer's response on them. At $\alpha=2$, equal logarithmic bands
carry equal first-moment energy, balancing that gain across
spectral scales.

\paragraph{Geometry seeking.}
We hypothesize that layers with $2<\alpha<3$ seek the geometry of
the data by amplifying structured variation across directions.
Their ideal mean is finite while their variance grows with the
upper cutoff. Higher bands contribute less to
the mean but more to the second moment. At $\alpha=2.5$, for
example, moving up one octave multiplies the first-moment band
weight by $1/\sqrt{2}$ and its second-moment weight by $\sqrt{2}$.

This spread measures how unevenly a layer amplifies its singular
directions. For the full
mean-normalized correlation operator $\widetilde{\mathbf X}$
with $M$ eigenvalues and $\langle\tilde\lambda\rangle=1$,
\[
 \operatorname{Var}_{\rm ESD}(\tilde\lambda)
   =\frac{1}{M}\sum_{i=1}^M(\tilde\lambda_i-1)^2
   =\frac{1}{M}\|\widetilde{\mathbf X}-I_M\|_F^2 .
\]
The normalized spectrum expresses these hypotheses through
concentration and directional spread. The fraction of the trace
carried by the leading modes measures concentration. The variance
of the eigenvalues measures the spread of squared directional
gains. In the proposed learning picture, mode seeking emphasizes
the first quantity and geometry seeking emphasizes the second.
\clearpage
